%! TEX root = rth.tex
% the main TeX file which is intended to compile, :VimtexReload after adjustment
   % this file should be included in the main file
% :h vimtex-tex-root
\documentclass[../rumukota.tex]{subfiles}
%ensure the class of docs

\begin{document}

\section{Ruang Terbuka Hijau}%
\label{sec:Ruang Terbuka Hijau}

\subsection{Pengertian Ruang Terbuka Hijau}%
\label{sub:Pengertian Ruang Terbuka Hijau}

Berikut versi yang telah diperbaiki dan tanpa sitasi:

Ruang Terbuka Hijau (RTH) adalah area berbentuk memanjang atau mengelompok yang penggunaannya bersifat terbuka dan menjadi tempat tumbuh tanaman, baik yang tumbuh secara alami maupun yang sengaja ditanam. RTH merupakan bagian dari ruang terbuka dalam penataan ruang perkotaan yang diisi oleh tumbuhan dan vegetasi, baik endemik maupun introduksi, yang memberikan manfaat langsung maupun tidak langsung, seperti meningkatkan keamanan, kenyamanan, kesejahteraan, dan keindahan lingkungan perkotaan.

Ruang terbuka dapat berbentuk jalur (path), seperti jalur hijau jalan, tepian waduk atau danau, bantaran sungai, sepanjang rel kereta api, serta koridor jaringan listrik tegangan tinggi. Selain itu, ruang terbuka juga dapat berbentuk simpul (nodes), seperti taman rumah, taman lingkungan, taman kota, taman pemakaman, maupun taman pertanian perkotaan.

Secara umum, ruang terbuka hijau dapat dipahami sebagai wilayah yang berbentuk memanjang atau mengelompok yang ditumbuhi tanaman dan memiliki berbagai manfaat bagi kehidupan. Keberadaan RTH diharapkan mampu menjadi penyeimbang lingkungan perkotaan, antara lain sebagai pengendali pencemaran udara, daerah resapan air, serta pengurang dampak polusi dari aktivitas kendaraan.

Berikut uraian lebih rinci mengenai jenis kepemilikan RTH:






\bibliographystyle{apalike}
\bibliography{../filename.bib} % required for citecompletion, not work in mainfile
\end{document}

